\documentclass[11pt]{article}
\usepackage[a4paper,margin=29mm]{geometry}
\usepackage{amsmath,amssymb,amsthm,mathtools,lmodern}
\usepackage[T1]{fontenc}
\DeclareMathSizes{10}{10}{8}{8}
\DeclareMathSizes{10.95}{10.95}{8}{8}
\usepackage[hidelinks]{hyperref}
\usepackage{microtype}
\setlength{\emergencystretch}{1.8em}
\newtheorem{proposition}{Proposition}[section]
\newtheorem{lemma}[proposition]{Lemma}
\newcommand{\wt}{\operatorname{wt}}
\newcommand{\pos}[1]{(#1)_+}
\title{The hereditary strict-peak reduction}
\hypersetup{unicode=true,
  pdftitle={The hereditary strict-peak reduction},
  pdfauthor={Zhang Yibo and Zhang Zhanhe}}
% The byline uses the title-page space otherwise reserved for an empty date.
\makeatletter
\renewcommand{\@maketitle}{%
  \begingroup\centering
  \vspace*{0.6em}%
  {\LARGE\@title\par}\vspace{0.85em}%
  \@author\par
  \endgroup
  \vspace{0.85em}}
\makeatother
\author{\begin{tabular}{@{}c@{\hspace{2em}}c@{}}
  \normalsize Zhang Yibo &
  \normalsize Zhang Zhanhe\\[-1pt]
  \footnotesize\href{https://orcid.org/0009-0002-2005-5781}{ORCID 0009-0002-2005-5781} &
  \footnotesize\href{https://orcid.org/0009-0007-2883-9640}{ORCID 0009-0007-2883-9640}
\end{tabular}\endgraf\vspace{3pt}%
{\footnotesize Both authors contributed equally as joint first authors. Zhang Zhanhe is the corresponding author.\endgraf}}
\date{}
\begin{document}
\maketitle
\begin{abstract}
This companion to \emph{Rank-unimodality of one-descent composition ideals}
gives the passage from a single increasing arm to the two interacting arms
at a peak. We specify the three coefficient comparisons, the finite
certificate in the mixed case, and the two inductions that produce the
ordinary rank sequence. The accompanying sources contain the proofs of
the coefficient estimates used in this reduction. The new transport
theorem in the article supplies its one unrestricted hypothesis.
\end{abstract}

\section{Notation and the two inductions}
Write \(S_t=z+\cdots+z^t\), \([t]=1+\cdots+z^{t-1}\),
\(\Delta v_n=v_n-v_{n-1}\), with zero extension. For an increasing arm
\(X\) let \(A=F_X\) be its ordinary rank polynomial, \(a=\wt X\) and
\(q_X=z^aA(1/z)\). The second arm \(Y\) has \(B,b,q_Y\) defined
in the same way. A legal append has capacity at least the previous maximum.
The identities \(F_{Zt}=1+S_tF_Z\) and \(q_{Zt}=[t]q_Z+z^{\wt(Zt)}\)
hold for such an append. An arm whose rank polynomial is an interval
is an elementary base. Every remaining arm can be normalized to
\(1^\alpha(t_1,\ldots,t_h)\) with \(1\le\alpha\) and
\(2\le t_1\le\cdots\le t_h\), preserving its weight and polynomial.

For \(P=AB\), let \(p=q_Xq_Y\) and \(u_n=\pos{\Delta p_n}\).
The count induction establishes unimodality of \(2^nu_n\). It deletes
one final non-one part in a heavier arm. The peak induction uses this
whole-sequence conclusion, with a proper relative of the same peak and
with cores at smaller peaks. The input to the count induction is the
following hypothesis at every final-part prefix of both arms. With
\(G_i=\sum_{v=0}^i A_v\), \(D_i=\pos{\Delta A_i}\) for \(i\ge2\),
it reads
\begin{equation}
 G_{k+1}D_n\ge G_kD_{n+1}
 \quad (0\le k\le n-3,\ D_nD_{n+1}>0).                    \label{C}
\end{equation}
Section 3 of the article proves \eqref{C} for each arm, with one
additional admissible read. For the following reductions, one can
equally well assume \eqref{C} on any final-deletion-closed class.

We use three established shape facts about the ordinary product. It is
unimodal; the positive differences of \(p\) have initial interval
support; and if both arms are noninterval arms then its first mode
\(\sigma\) is strictly less than \(\max(a,b)\), including \(a=b\).
These are statements about the ordinary product before inserting a peak.
The underlying interval-layer, midpoint and positive-difference proofs
are collected in the evidence index in Section~\ref{index}.

Here is the midpoint argument. The lowest height layer of \(A\) is
\([0,a]\), and every higher interval has endpoint sum at least \(a+1\).
Appending \(t\) sends \(A\) to \(1+S_tA\); interval convolution gives
the new higher-layer bound, and the constant changes the lowest layer.
For a noninterval arm its second height layer is exactly \([2,a-1]\).
Thus the layer products of \(A,B\) give growth through
\(\lceil(a+b)/2\rceil\). For \(a=b\), the two lowest layers contribute
\(-1\) on the edge from \(a\) to \(a+1\). The two second layers together
contribute \(+1\), and all other pairings contribute nonnegatively.
Orient \(a\ge b\). In \(AB=B+S_tF_ZB\) the second term is unimodal by
deletion induction and interval convolution. When \(a>b\) the known
growth extends past \(b\), and when \(a=b\) the extra edge gives the
same conclusion. Thereafter only the unimodal window remains.
This proves ordinary unimodality and the strict reverse mode bound.
The positive-difference interval support is proved, also by window
and deletion, in \path{joint_rising_differences.tex}.

\section{When a sliding interval preserves weighted differences}
Let \(m\) be the last ordinary mode of \(p\) and \(\ell\) the first mode
of \(2^nu_n\). Define \(E_i=2p_i-p_{i-1}\).
We need ordinary unimodality of the tilted coefficients \(2^np_n\) and
the following local implication for \(t\ge1,\ n>m,\ k=n-t\ge0,\ k<\ell\):
\begin{equation}
 \left.
 \begin{gathered}
 p_n>p_k,\quad p_{n+1}>p_{k+1},\quad 0\le E_n\le E_k
 \end{gathered}\right\}
 \quad\Longrightarrow\quad
 E_{n+1}-E_n\le E_{k+1}-E_k.                               \label{gap}
\end{equation}
Equality at \(E_n=E_k\) is included.

\begin{lemma}\label{win}
If \(2^np_n\) and \(2^nu_n\) are unimodal and \eqref{gap} holds,
then the weighted positive differences of every \([t]p\) are unimodal.
\end{lemma}
\begin{proof}
Put \(w_n=\Delta([t]p)_n=p_n-p_{n-t}\), \(f_n=2^n(w_n)_+\).
The support of the positive differences is initial: while the outgoing
index precedes the first mode and the entering index passes the last
mode, their difference decreases; once the outgoing index passes the
mode it is nonpositive. Before that crossing it is positive up to a
possible terminal plateau, by the initial positive support for \(p\).
For \(n\le m\), \(f_n=2^n([t]u)_n\), which agrees with convolution of
the unimodal \(2^nu_n\) against \((1,2,\ldots,2^{t-1})\).
Such a log-concave interval kernel preserves unimodality: the signed
first differences of the input have one change of sign, and the
nonnegative order-two minors of the kernel preserve this sign order.
At \(m+1\), clipping the signed sum gives at most that comparison
sequence. Alternatively, for \(n<t\) use the unimodal comparison
\(2^np_n=f_n\) with \(f_t\le2^tp_t\). Choose the appropriate comparison
through \(c=\max(m,t-1)\). Any earlier strict fall forces
\(f_{c+1}\le f_c\). This checks the first crossing boundary.

Consider now a nonstrict fall at \(n\ge c+1\) while
\(w_n,w_{n+1}>0\). If \(k=n-t\ge\ell\), then
\(w_n\le\sum_{i>k}u_i\le2u_{k+1}\), by the successive halvings of
the finite tail. Positivity of \(w_n\) puts \(k\) before the first
ordinary mode. Consequently
\(
 w_{n+1}=w_n-u_{k+1}-(p_n-p_{n+1})\le w_n/2.
\)
If \(k<\ell\), the fall gives \(E_n-E_k=2w_n-w_{n-1}\le0\).
Here \(E_{k+1}>0\) and \(E_{k+1}-E_k=2u_{k+1}-u_k\ge0\).
For \(E_n<0\), tilted unimodality implies \(E_{n+1}\le0\).
For \(E_n\ge0\), \eqref{gap} gives \(E_{n+1}-E_{k+1}\le E_n-E_k\).
In both cases \(2w_{n+1}-w_n\le0\). A nonstrict fall propagates
at every remaining positive edge, and all terms after the support are
zero. This excludes a strict fall followed by a rise.
\end{proof}

\section{The coefficient comparisons for genuine increasing arms}
On the forward side put \(E_Y^*=1-z+(2z-1)B\),
\(L=E_Y^*A\), and \(e_i=2P_{i-1}-P_i\).
When both factors are strictly rising through \(n+1\),
the source increments \(w_j=\Delta E_Y^*(j)\) are nonnegative through
that rank and vanish for \(j<3\). The single-arm proof obtains this from
the ordering of the ordinary and source modes. Summation by parts gives
\(L_i=\sum_{j=3}^iw_jG_{i-j}\) at the two ranks used. Hypothesis
\eqref{C} gives \(L_{n+1}D_n\ge L_nD_{n+1}\), and \(L_{n+1}\ge L_n\).
Since \(e_i=L_i-D_i\) here, \(e_{n+1}<0\) implies \(e_n\le0\);
if \(e_{n+1}\ge0\), then
\[
 L_{n+1}(e_{n+1}-e_n)
 = L_{n+1}D_n-L_nD_{n+1}+(L_{n+1}-L_n)e_{n+1}\ge0.
\]
Outside the common rising prefix, the first nonincreasing factor can
be oriented as \(A\). The identity \(e_i=L_i-\Delta A_i\) and
nonnegativity of the whole \(E_Y^*\) give \(e_i\ge0\) at those ranks.
These comparisons supply unimodality of \(2^np_n\) after reversal and
the common-rising part of \eqref{gap}; the initial ranks are included
in the source proof. A negative current reverse tilt has no positive
successor.

For the remaining branches write \(d=a+b,\ j=d-n\), and
\(T=E_{n+1}-E_n=P_{j+1}-3P_j+2P_{j-1}\).
If both forward factors have stopped rising, the first-half
both-nonincreasing theorem gives \(T\le0\) on its forward half.
The forward condition is \(2j\le d+2\). On the remaining ranks the ordinary reverse half-tail estimate gives
\(T\le0\) for \(n>m,\ 2n\le d\), without a sign condition on \(E_n\).
Together they cover the comparison in \eqref{gap}, since the earlier
increment there is nonnegative. The proofs keep the other half and the
source correction in distinct rank ranges.

It remains to compare the mixed phase. Orient \(X=1^\alpha(t_1,\ldots,t_h)\)
as the forward nonincreasing owner, with rising opposite \(Y\).
Put \(r=a-j\), \(s_r=\Delta q_X(r)\), and \(C_0=s_{r+1}-s_r\).
The exact signed identity is
\begin{equation}
 T=C_0-\sum_i \Delta E_Y^*(i)A_{j+1-i}.                         \label{mixed}
\end{equation}
Let \(K\) be the first weight at which some composition fails to fit
\(Y\); it is the first index of the positive source \(E_Y^*\).
For \(Y=1^\beta(v_1,\ldots,v_g)\) noninterval this equals
\(\min(\{g+2\}\cup\{v_i+g-i+1:1\le i\le g\})\).
A missing one-part word or a shortest failure at the end of a suffix
gives the formula; equivalently it follows recursively under a legal
append. The witnesses concern distinct words in generalized subword order.

Here are the four pieces establishing \eqref{gap} in the mixed phase.
When both first missing weights are at most twelve the finite
reduction below and the through-nine base already give the count step.
Otherwise an opposite with \(K\le9\) has an owner of first missing
weight at least thirteen and hence at least eleven non-one parts.
The through-eight and through-nine source estimates exclude \(T>0\)
for those owners. A one-non-one-part opposite has \(K=3\); the
first-source reflection proof includes it. At \(r\le-2\) there is
no positive curvature; an unpaid \(r=-1\) birth has size one and
\(j\le K-2\). It falls in the strict opposite prefix in each branch.

For \(10\le K\le12\), the exact initial exchange permits
\(t_1-1\le\alpha\). At any remaining positive-curvature rank with
\(r\ge0\), put \(d_0=K-\alpha\) and \(L=\sum(t_i-1)=2r+2+e\).
Unbounded estimates give
\[
 d_0\ge6,\qquad 0\le e\le d_0-3,\qquad
 L\le\min(2d_0(d_0-3),12d_0-25).
\]
For \(d_0=6\) only \(h\le3\) remains and for \(7\le d_0\le8\)
only \(h\le5\) remains. With \(S_K=\lfloor(1+(K-2)(K-1)/2)/2\rfloor\),
the possible ranks partition as follows. The final column requires
enumerating every sorted partition of \(L\) allowed by the height
bounds in the archived certificate.
\[
\begin{array}{c|r|rrr}
 K&\text{rank states}&\text{strict prefix}&\text{source ratio}&
     \text{remaining}\\ \hline
 10&684&170&451&63\\ 11&1048&279&641&128\\ 12&1509&410&891&208
\end{array}
\]
Of the resulting \(1{,}858{,}242\) owner/rank inputs,
\(72{,}043\) require the final positive-curvature payment. Each
satisfies the exact integer inequality
\[
 C_0\le A_{j+1-K}+\sum_{i=2}^{K-3}2^{i-2}A_{j+1-K-i}.
\]
This is a finite part of the proof, with an exhaustive domain derived
from the preceding inequalities. Thus \(T>0\) implies \(j-2\le S_K\).
The two positive windows imply \(k+1\le j-2\). The strict prefix
has \(\Delta q_Y(s)\ge1\) for \(0\le s\le S_K\), by the interval
products and their centers. When \(0\le r<k\), the \(r+1\) term in
\([z^{k+1}](2-z)q_X(1-z)q_Y\) alone pays \(C_0\). The constant
term pays \(r=-1\). When \(r\ge k\) the relative source comparison
gives \(E_{k+1}-E_k\ge\lambda E_k+1\), where
\(E_n\ge d_r>0,\ \lambda=C_0/d_r\). This proves \eqref{gap} also at equality.

For \(K=13,14,15\), let \(M_Y=\beta+\sum(v_i-1)\),
\(B_Y=\max(\beta,\max(v_i-1))\), and
\(s_Y=\min(\lfloor M_Y/2\rfloor,M_Y-B_Y)\).
The last eleven non-one parts dominate \(2,\ldots,12\).
Expanding the last intervals gives the uniform bound
\(20\,\Delta B_i\le2^i\) for every \(i>s_Y\).
The expansion retains all seven proper suffix terms. In the bulk
their combined coefficient is at most \(i^5\); the long kernel gives
the remaining density at most \(1/25\) and the suffixes at most
\(1/100\). Outside the bulk the largest interval gives a decreasing
geometric bound.
For a positive curvature with \(j-2>s_Y\), the three integer
guards imply \(e\le2r+2\) and \(h+r+1\le5(r+1)\).
If \(Q=q_X(r+1)\), \(R=2q_X(r)-q_X(r-1)\), then \(Q\le5R\);
here \(C_0=Q-R\).
The source sum in \eqref{mixed} is at least
\(Q(1-4\,2^{-(j+1)}\Delta B_{j+1})\ge4Q/5\ge C_0\).
All remaining recoveries therefore lie in the strict prefix, where
the same earlier increment pays them. For \(K\ge16\) the corresponding
fourteen-terminal-part estimate \(i\,\Delta B_i\le2^i\) for \(i>s_Y\)
combines with the indexed owner bound. It puts every positive
recovery in the strict prefix. The details and exact constants for
both estimates are in the indexed source proofs.
This exhausts the mixed case, with the old weighted difference shape
used only to make the earlier increment nonnegative or to invoke
Lemma~\ref{win}.

\section{Count induction and the peak join}
If an arm is an interval, the rectangular base applies with arbitrary
opposite. The heavier-one-non-one-part base also allows an arbitrary
opposite. Otherwise remove \(t\) from the heavier new arm \(X=Zt\).
By induction \(q_Zq_Y\) has unimodal weighted positive differences.
The three preceding comparisons and Lemma~\ref{win} give the same for
\([t]q_Zq_Y\). At this stage both old arms are noninterval, since the heavier
one-non-one-part case was already handled. Now
\(
 q_Xq_Y=[t]q_Zq_Y+z^a q_Y.
\)
Its first mode is strictly below \(a\). Through that mode the positive
differences are the window differences, and thereafter the ordinary
product is nonincreasing. Truncating the unimodal weighted positive
differences at that index completes the count induction.

For a strict peak \(M>\max(X,Y)\), put
\[
 H_M=A+B+(S_M-1)P,\qquad
 Q_{X,Y}=A+(S_M-1)P,\qquad Q_{Y,X}=B+(S_M-1)P .
\]
Set \(\delta_M(Y)=([M]-z^M)q_Y+z^{b+M}\).
The reversal of \(Q_{X,Y}\) about \(a+b+M\) is \(q_X\delta_M(Y)\).
Choose a least failing peak and then a failing relative with the
fewest total arm parts at that peak. The balanced-deletion and outer
layer lemmas eliminate interval arms and allow the orientation
\(a\le b\), \(b-a\le M-4\). Within this band either relative is
unimodal exactly when the core is. The balance and equivalence follow
from the deletion and outer-layer lemmas in the integrated foundation.

Put \(N=a+b+M\), \(R_0=\lfloor N/2\rfloor\),
\(W=([M]-z^M)p\). Through \(R_0\) the reverse core agrees with \(W\):
its two corrections start at \(a+M,b+M>R_0\). Write \(Y=Zt\).
The unimodal proper relative \(q_X\delta_M(Z)\), after convolution
with \([t]\), agrees with \(W\) through \(b-1\). Indeed the corrections
in this identity have factors \(z^b\) and \(z^{b-t+M}\).
Consequently \(\Delta W\) cannot go from negative to positive on that
prefix. For \(0\le n<M\) the identity
\((1-z)W=(1-2z^M+z^{M+1})p\) gives \(\Delta W_n=p_n\ge0\).
The two initial ranges therefore cover all \(n<\max(b,M)\), including
the boundary \(b=M\).

At \(\max(b,M)\le n=M+k\le\min(R_0,M+\sigma)\),
\[
 \Delta W_{M+k}=p_{M+k}-p_k-u_k .
\]
If \(k<\ell\), then \(2u_{k+1}\ge u_k\), hence \(p_k+u_k\)
is nondecreasing and \(p_{M+k}\) is nonincreasing since \(b>\sigma\).
When \(b>M\) this comparison includes the joining edge from \(b-1\),
since \(b-1\ge\sigma\). When \(b\le M\) all earlier differences through
\(M-1\) are nonnegative. If \(\ell\le k\),
then \(p_{M+k}-p_k\le\sum_{i>k}u_i\le u_k\) by the weighted
halvings. This includes \(k=\sigma\). Thus the signs have the required
order through \(M+\sigma\) or the end of the relevant half.

For \(M+\sigma<n\le R_0\), write \(s=a+b-\sigma\). The range is
nonempty only if \(s>\max(a,b)\). For the largest arm part \(t_0\),
the base \(H_{t_0}\) is unimodal: strip all end copies of \(t_0\),
apply the strictly smaller peak case (or the interval base), and
restore each copy by \(F\mapsto1+S_{t_0}F\).
The last forward mode of \(P\) is \(s\), so
\[
 \Delta H_{t_0}(s+1)=2P_s-P_{s+1}-P_{s-t_0}>0.
\]
Its unimodality gives forward growth through \(s+1\). As
\(H_M=H_{t_0}+z^{t_0+1}[M-t_0]P\), at each corresponding \(j=N-n<s\)
we have \(\Delta H_M(j+1)=\Delta H_{t_0}(j+1)+P_{j-t_0}-P_{j-M}\ge0\).
Forward first-half growth covers the rest. This contradicts balanced
core-relative equivalence. In particular both relatives and the core
are unimodal for every strict peak. The factorization in the article
restores repeated maximum letters and the final increasing suffix.

\section{Evidence index}\label{index}
The supplement contains the following source families. The numbered
proofs in \emph{Detailed coefficient lemmas} typeset the decisive
reduction fragments. Their original files and the data used by the finite
certificates are under \path{reduction/}. The foundation, including the
full proof of the balanced reduction, is under \path{foundations/}.
Historic source comments about the remaining hypothesis refer to the
stage before the transport theorem of the article.
\begingroup\small\raggedright\def\UrlBreaks{\do\_\do\.\do\/}\begin{description}
\item[Ordinary shape and deletion.] The integrated foundation's
sections on structure and minimal counterexamples, plus
\path{formal_arm_extensions.tex} and \path{joint_rising_differences.tex}.
\item[Source signs.] \path{clipped_forward_ratio_frontier.tex},
\path{C_to_coupled_S_rising_window.tex},
\path{agent_bothdeclining_first_birth.tex}, and
\path{agent_half_lc_actual_weak_curvature.tex}.
\item[Closed interval comparison and certificates.]
\path{coupled_source_window_A2_and_kappa_nine_original_proof.tex},
\path{first_missing_twelve_original_criterion_proof.tex},
\path{fixed_terminal_fourteen_slope_thinness_proof.tex}, and
\path{complete_mixed_source_unrestricted_C_original_criterion_proof.tex}.
The exact finite domain through twelve is stated in equations
(1)--(7) of the corresponding source and checked by its certificate.
\item[Two inductions.] \path{one_hereditary_C_arm_arbitrary_other_original_proof.tex},
\path{inductive_signed_window.tex}, \path{short_prefix_reflection_join.tex}
and \path{general_base_join_route.tex}. The rectangular and heavier-one-part
bases are in \path{binary_ternary_A2_original_family_proof.tex} and
\path{deficit_cumulative_ratios.tex}.
\item[New formal and enumerative checks.] \path{DepthSlackTransport.lean}
contains the finite mass transport, the induction over all budgets and
the delayed ratio transfers. \path{DepthSlackArmIdentities.lean} checks
the exact polynomial identities for every finite list of capacities. The two independent enumerations in
\path{checks/} verify actual ideals and selected hypotheses directly.
\end{description}\endgroup
\end{document}